\chapter{多电子原子与谱项符号}\label{ch:b3:many-electron-atoms}

1868 年，太阳光谱中的一条黄线揭示了一种地球上未知的元素，
即氦。1895 年，它终于在实验室中被分离出来，
其光谱却像是属于\emph{两}种元素：两族谱线，
一度被归于“正氦”和“仲氦”，
且两族谱线之间从不互相发射或吸收光。其实
只有一种氦。这两族谱线分别是同一原子的三重态与单重态，
它们由电子的自旋，以及一条比任何力都更基本的规则而分开：
两个电子交换时，波函数必须变号。
本章沿着这条规则，从泡利原理出发，
一直推导到光谱学家用来标记每一种原子的每一个状态的谱项符号，
以及每一种离子——包括过渡金属离子，其颜色
在\cref{ch:b3:complex-spectra-magnetism}中得到解释。

\begin{recall}
本科一年级分册用四个量子数 $n$、$l$、
$m_l$、$m_s$ 标记电子，并借助泡利原理、克莱奇科夫斯基规则
与洪特规则构造基态电子排布，并数出未成对电子。本科二年级分册
给出了轨道能量，以及用于计算屏蔽的斯莱特规则。
\Cref{ch:b3:quantum-model-systems}给出了算符、本征值，以及转子的角动量
（$\hat L^2$ 的本征值为 $l(l+1)\hbar^2$）和
氢原子。
\end{recall}

\section{自旋与自旋轨道}

电子带有一种内禀角动量，即自旋，其量子数为
$s = \frac12$：$\hat S^2$ 只有一个本征值 $s(s+1)\hbar^2 =
\frac34\hbar^2$，$\hat S_z$ 有两个本征值 $m_s\hbar = \pm\frac12\hbar$。
两个自旋函数分别记为 $\alpha$（$m_s = +\frac12$）和 $\beta$
（$m_s = -\frac12$），它们是正交归一的：$\langle\alpha|\alpha\rangle =
\langle\beta|\beta\rangle = 1$，$\langle\alpha|\beta\rangle = 0$。自旋没有经典对应，
它不是小球的自转，而且
在非相对论哈密顿量中根本不出现。

\begin{definition}[自旋轨道]\label{def:b3:many-electron-atoms:spin-orbital}
\emph{自旋轨道}\index{自旋轨道}是空间轨道
$\phi(\mathbf r)$ 与自旋函数的乘积，即 $\phi\alpha$ 或 $\phi\beta$。
自旋轨道给出了一个电子的完整描述。
\end{definition}

本科一年级分册中用 $(n, l, m_l)$ 标记的原子轨道，
由此得到两个自旋轨道，这就是一个电子亚层可容纳 $2(2l+1)$ 个电子的原因。
真正的问题在于，为什么一个自旋轨道至多只能容纳一个电子。

\section{不可区分电子与斯莱特行列式}

\begin{definition}[不可区分粒子，反对称波函数]\label{def:b3:many-electron-atoms:indistinguishable}
若任何测量都无法区分两个粒子，则称它们为\emph{不可区分的}\index{不可区分粒子}。
即每个电子都与其他每个电子完全相同。若多电子波函数在交换
任意两个电子的坐标（空间与自旋）时变号，则称它为
\emph{反对称波函数}\index{反对称波函数}，其满足
下式：
$\Psi(\dots,i,\dots,j,\dots) = -\Psi(\dots,j,\dots,i,\dots)$。
\end{definition}

仅凭不可区分性，只能要求交换前后 $|\Psi|^2$ 保持不变，
因而 $\Psi$ 至多只能变号。自然界的选择是：
任何电子体系的波函数——也就是任何费米子体系的波函数——都是
反对称的。这是量子力学的一条公设，
每一条原子光谱都证实了这一点。

\begin{definition}[斯莱特行列式]\label{def:b3:many-electron-atoms:slater}
对于 $N$ 个电子占据 $N$ 个不同的自旋轨道 $\chi_1, \dots, \chi_N$，
\emph{斯莱特行列式}\index{斯莱特行列式}为
\[
\Psi = \frac{1}{\sqrt{N!}}\begin{vmatrix}
\chi_1(1) & \chi_2(1) & \cdots & \chi_N(1)\\
\chi_1(2) & \chi_2(2) & \cdots & \chi_N(2)\\
\vdots & & & \vdots\\
\chi_1(N) & \chi_2(N) & \cdots & \chi_N(N)
\end{vmatrix},
\]
其中第 $i$ 行给出电子 $i$ 在各自旋轨道中的取值，简记为
$|\chi_1\chi_2\dots\chi_N|$。
\end{definition}

\begin{theorem}[泡利原理]\label{thm:b3:many-electron-atoms:pauli}
斯莱特行列式是反对称的，若两个电子占据同一自旋轨道，
则该行列式为零。因此，原子中任何两个电子都不可能具有
完全相同的四个量子数。
\end{theorem}

\begin{proof}
交换电子 $i$ 与 $j$ 相当于交换行列式的两行，
行列式因而变号，对任意一对电子反对称性都成立。若两个
自旋轨道相同，则两列相等，行列式为零：
这样的状态不存在。同一原子中两个电子若具有相同的 $n$、$l$、
$m_l$ 与 $m_s$，就会占据同一自旋轨道。
\end{proof}

\begin{example}[氦的基态]\label{ex:b3:many-electron-atoms:helium-ground}
两个电子都处于 $1s$ 轨道时，
\[
\Psi = \frac{1}{\sqrt2}\begin{vmatrix}1s\alpha(1) & 1s\beta(1)\\ 1s\alpha(2) &
1s\beta(2)\end{vmatrix} = 1s(1)\,1s(2)\cdot\frac{1}{\sqrt2}\big[\alpha(1)\beta(2)
- \beta(1)\alpha(2)\big].
\]
空间部分是对称的，自旋部分是反对称的：两个自旋
配对，总自旋为零。
\end{example}

\section{氦：库仑排斥与交换}

将氦的一个电子激发到 $2s$。
由 $1s\alpha$、$1s\beta$、$2s\alpha$、$2s\beta$ 出发，每个轨道各放一个电子，可构造出四个行列式。
重新组合后，它们可写成一个空间函数与一个自旋函数的乘积：
\[
\Psi_\pm = \frac{1}{\sqrt2}\big[1s(1)2s(2) \pm 2s(1)1s(2)\big]\times(\text{自旋
函数}).
\]
反对称的空间函数 $\Psi_-$ 必须与一个对称的
自旋函数配对。对称的自旋函数共有三个：$\alpha(1)\alpha(2)$、
$\beta(1)\beta(2)$ 以及 $\frac{1}{\sqrt2}[\alpha(1)\beta(2) + \beta(1)\alpha(2)]$，
这三者为 $M_S = 1, 0, -1$ 三个分量，对应总自旋 $S = 1$。
对称的 $\Psi_+$ 必须与唯一的反对称自旋函数配对，此时 $S = 0$。

\begin{definition}[单重态与三重态]\label{def:b3:many-electron-atoms:singlet-triplet}
总自旋 $S = 0$ 的态称为\emph{单重态}\index{单重态}；
总自旋 $S = 1$ 的态，其三个分量 $M_S = 1, 0, -1$ 在无磁场时
能量相同，称为\emph{三重
态}\index{三重态}。
\end{definition}

\begin{definition}[交换积分]\label{def:b3:many-electron-atoms:exchange}
设 $a$、$b$ 为两个轨道，电子排斥取 $e^2/4\pi\varepsilon_0
r_{12}$，则\emph{交换积分}\index{交换积分}为
\[
K_{ab} = \iint a(1)b(2)\,\frac{e^2}{4\pi\varepsilon_0r_{12}}\,b(1)a(2)\,\dd\tau_1\dd\tau_2,
\]
它与经典电子排斥积分 $J_{ab}$（两边均为
$a(1)b(2)$）形式相同，区别仅在于一侧的电子互换了。
$K_{ab}$ 为正，且没有经典对应量。
\end{definition}

\begin{proposition}[单重态与三重态的能量]\label{prop:b3:many-electron-atoms:helium}
对电子排斥作一级近似时，氦的 $1s2s$ 态
能量为 $E_\pm = E_{1s} + E_{2s} + J \pm K$，单重态（$+$）
比三重态（$-$）高 $2K$。
\end{proposition}

\begin{proof}
$\Psi_\pm$ 的未微扰能量均为 $E_{1s} + E_{2s}$（核电荷为 2 的类氢
轨道）。一级修正是
归一化态中排斥算符 $\hat V_{12}$ 的期望值：
$\frac12\langle 1s(1)2s(2) \pm 2s(1)1s(2)|\hat V_{12}|1s(1)2s(2) \pm
2s(1)1s(2)\rangle$。两个直接项给出 $\frac12(J + J)$；两个交叉
项给出 $\pm\frac12(K + K)$，因为 $\hat V_{12}$ 对
电子是对称的。故 $\Delta E = J \pm K$，自旋函数已归一化
且不受 $\hat V_{12}$ 影响，仅乘以 1。
\end{proof}

三重态的能量反而更低，尽管哈密顿量中不含自旋。其
空间函数在 $\mathbf r_1 = \mathbf r_2$ 时为零，因此两个
电子彼此回避，排斥较弱。这一“费米空穴”正是
洪特第一规则的物理内涵。

\begin{example}[氦的交换积分：实测值]\label{ex:b3:many-electron-atoms:K}
% ledger: lev:He.2s3S1, lev:He.2s1S0
氦的 $1s2s$ 能级在基态之上分别位于 \qty{159856}{cm^{-1}}（${}^3$S）和
\qty{166277}{cm^{-1}}（${}^1$S）处：单重态高出
\qty{6421}{cm^{-1}}，即 \qty{0.796}{eV}，且 $K(1s,2s) = \qty{0.398}{eV}$。
下图给出两族能级。
\end{example}

\section{罗素–桑德斯耦合与谱项符号}

在轻原子中，电子的轨道角动量相加得到
总轨道角动量 $\mathbf L$，自旋相加得到总自旋 $\mathbf S$，然后
$\mathbf L$ 与 $\mathbf S$ 才发生较弱的耦合。

\begin{definition}[罗素–桑德斯耦合]\label{def:b3:many-electron-atoms:russell-saunders}
在\emph{罗素–桑德斯耦合}\index{罗素–桑德斯耦合}中，一个组态的各态
由\emph{总轨道角
量子数}\index{总轨道角量子数}
$L$（$\hat{\mathbf L}^2 = L(L+1)\hbar^2$，$M_L = \sum m_l$）、\emph{总自旋
量子数}\index{总自旋量子数} $S$（$M_S = \sum m_s$）、
以及\emph{总角量子数}\index{总角量子数}
$J$，其取值为 $|L - S|, \dots, L + S$。
\end{definition}

\begin{definition}[光谱项、自旋多重度、谱项符号]\label{def:b3:many-electron-atoms:term}
\emph{光谱项}\index{光谱项}是某一组态中
$(2L+1)(2S+1)$ 个态的集合，这些态具有相同的 $L$ 与 $S$。其
\emph{自旋多重度}\index{自旋多重度}为 $2S + 1$。\emph{谱项符号}
\index{谱项符号}为 ${}^{2S+1}L$，字母 S、P、D、F、G、H
依次对应 $L = 0, 1, 2, 3, 4, 5$；给定 $J$ 的态写作 ${}^{2S+1}L_J$ ——
例如 \termsym{3}{P}{2}。
\end{definition}

\begin{proposition}[闭壳层]\label{prop:b3:many-electron-atoms:closed-shell}
满电子亚层有 $L = 0$ 且 $S = 0$，
对谱项符号没有贡献。
\end{proposition}

\begin{proof}
在满电子亚层中，每个 $m_l$ 都出现两次，每个 $m_s = \pm\frac12$
都出现 $2l+1$ 次，故 $M_L = 0$，$M_S = 0$；且只有一种
填充方式（一个行列式）。唯一的 $M_L = M_S = 0$ 的态
只能属于 $L = 0$，$S = 0$。
\end{proof}

\begin{proposition}[组态态数的计数]\label{prop:b3:many-electron-atoms:count}
$N$ 个电子占据 $2(2l+1)$ 个自旋轨道的亚层，给出
$\binom{2(2l+1)}{N}$ 个行列式，且其各项的简并度 $(2L+1)(2S+1)$
之和等于此数。
\end{proposition}

\begin{proof}
行列式由被占据的自旋轨道集合唯一确定（其排列次序
只改变符号）：即选 $N$ 个，取自 $2(2l+1)$ 个。各项即为
同一批态的重新分组，故计数必须一致。
\end{proof}

\begin{method}[组态的谱项]\label{met:b3:many-electron-atoms:terms}
\begin{enumerate}
  \item 列出满足泡利原理的全部行列式，
        并按 $M_L$ 与 $M_S$ 制表。
  \item 找出最大的 $M_L$；与之相伴的最大 $M_S$
        即开始一个谱项，$L = M_L$，$S = M_S$。
  \item 对每个数对 $(M_L, M_S)$（满足
        $|M_L| \le L$、$|M_S| \le S$），划去一个行列式。
  \item 对剩余部分重复上述步骤，直至表格为空；核对计数。
\end{enumerate}
\end{method}

\begin{example}[$p^2$ 的谱项]\label{ex:b3:many-electron-atoms:p2}
两个处于 $p$ 轨道（$m_l = 1, 0, -1$）的电子给出 $\binom62 = 15$ 个行列式。
按 $M_L$ 与 $M_S$ 列表如下：
\begin{center}\small
\begin{tabular}{c|ccc}
\hline
 & $M_S = 1$ & $M_S = 0$ & $M_S = -1$\\ \hline
$M_L = 2$ & -- & $(1^+,1^-)$ & --\\
$M_L = 1$ & $(1^+,0^+)$ & $(1^+,0^-)$, $(1^-,0^+)$ & $(1^-,0^-)$\\
$M_L = 0$ & $(1^+,-1^+)$ & $(1^+,-1^-)$, $(1^-,-1^+)$, $(0^+,0^-)$ & $(1^-,-1^-)$\\
$M_L = -1$ & $(0^+,-1^+)$ & $(0^+,-1^-)$, $(0^-,-1^+)$ & $(0^-,-1^-)$\\
$M_L = -2$ & -- & $(-1^+,-1^-)$ & --\\ \hline
\end{tabular}
\end{center}
（其中 $m_l^\pm$ 表示 $m_s = \pm\frac12$）。$M_L = 2$ 只与 $M_S = 0$ 同时出现：
得一个 ${}^1$D 项（5 个态）。剩下的最大 $M_L = 1$ 与
$M_S = 1$ 相伴：得一个 ${}^3$P 项（9 个态）。尚余一个 $M_L = M_S = 0$ 的态：
即一个 ${}^1$S 项。确有 $5 + 9 + 1 = 15$。
\end{example}

\begin{example}[碳，实测]\label{ex:b3:many-electron-atoms:carbon}
% ledger: lev:C.3P1, lev:C.3P2, lev:C.1D2, lev:C.1S0
碳的基态组态 $2p^2$ 给出能级 \termsym{3}{P}{0}，
\termsym{3}{P}{1}、\termsym{3}{P}{2} 分别位于 0、16.4 和 \qty{43.4}{cm^{-1}}，随后
\termsym{1}{D}{2} 位于 \qty{10193}{cm^{-1}}，\termsym{1}{S}{0} 位于
\qty{21648}{cm^{-1}}：${}^3$P 项最低，与洪特规则的预言相符。
\end{example}

\begin{omfigure}
\begin{tikzpicture}
\begin{axis}[name=main, width=0.62\linewidth, height=6.4cm, xmin=-0.3, xmax=6.6,
  ymin=-1500, ymax=23500, axis y line=left, axis x line=none,
  ylabel={能量 / cm$^{-1}$}, unbounded coords=jump, clip=false,
  scaled y ticks=false, ytick={0,10000,20000},
  tick label style={font=\small}, label style={font=\small}]
\addplot[omstate] table[x=x, y=E] {figdata/out/bachelor-3-atomic-levels-carbon.dat};
\draw[densely dotted, gray] (axis cs:1,4858.5) -- (axis cs:2,29.6);
\draw[densely dotted, gray] (axis cs:1,4858.5) -- (axis cs:2,10192.7);
\draw[densely dotted, gray] (axis cs:1,4858.5) -- (axis cs:2,21648);
\draw[densely dotted, gray] (axis cs:3,10192.7) -- (axis cs:4,10192.7);
\draw[densely dotted, gray] (axis cs:3,21648) -- (axis cs:4,21648);
\node[font=\small, anchor=south] at (axis cs:0.5,4858.5) {$2p^2$};
\node[font=\small, anchor=south] at (axis cs:2.5,21648) {${}^1$S};
\node[font=\small, anchor=south] at (axis cs:2.5,10192.7) {${}^1$D};
\node[font=\small, anchor=south] at (axis cs:2.5,29.6) {${}^3$P};
\node[font=\small, anchor=west] at (axis cs:5.05,21648) {\termsym{1}{S}{0}};
\node[font=\small, anchor=west] at (axis cs:5.05,10192.7) {\termsym{1}{D}{2}};
\node[font=\scriptsize, anchor=north] at (axis cs:0.5,-600) {组态};
\node[font=\scriptsize, anchor=north] at (axis cs:2.5,-600) {项};
\node[font=\scriptsize, anchor=north] at (axis cs:4.5,-600) {能级};
\draw[gray, ->] (axis cs:3.05,60) -- (axis cs:6.9,2500);
\end{axis}
\begin{axis}[at={(main.south east)}, anchor=south west, xshift=1.2cm, yshift=0.9cm,
  width=3.6cm, height=3.6cm, xmin=-0.1, xmax=1.1, ymin=-5, ymax=50,
  axis y line=right, axis x line=none, unbounded coords=jump, ytick={0,16.4,43.4},
  yticklabels={0,16.4,43.4}, tick label style={font=\scriptsize}, clip=false,
  title={\scriptsize ${}^3$P，放大图}, title style={yshift=-4pt}]
\addplot[omstate] table[x=x, y=E] {figdata/out/bachelor-3-atomic-levels-carbon-zoom.dat};
\node[font=\scriptsize, anchor=east] at (axis cs:0,0) {$J = 0$};
\node[font=\scriptsize, anchor=east] at (axis cs:0,16.4) {$J = 1$};
\node[font=\scriptsize, anchor=east] at (axis cs:0,43.4) {$J = 2$};
\end{axis}
\end{tikzpicture}

{\small 碳的基态组态：一个组态，三个项
（各项取实测能量；${}^3$P 项取其各能级的加权平均，
组态取全部十五个态的平均值），以及五个
能级。${}^3$P 项的自旋–轨道分裂（放大图，单位为
\unit{cm^{-1}}）比项与项之间的
间距小几百倍。}
\end{omfigure}

\section{洪特规则与自旋–轨道耦合}

\begin{proposition}[项的洪特规则]\label{prop:b3:many-electron-atoms:hund-terms}
对原子或离子的基态组态，能量最低的项是
满足（1）$S$ 最大，再在该 $S$ 下（2）$L$ 最大的那一项。其最低
能级满足（3）：不足半满时 $J = |L - S|$，超过半满时 $J = L +
S$。
\end{proposition}
\begin{proof}[地位]
这些规则由实验确立，仅适用于基态组态；
它们并非定理。规则 1 即上一
节所述的交换稳定化；规则 2 指自旋同向的电子
相遇较少；规则 3 由下文自旋–轨道耦合的符号得出。
\end{proof}

\begin{method}[由轨道框图求基态项]\label{met:b3:many-electron-atoms:ground-term}
\begin{enumerate}
  \item 自 $m_l = +l$ 起向下填充开壳层的各轨道框，
        每框先填一个自旋平行的电子，再从 $+l$ 起配对。
  \item $S = \frac12 \times$（未成对电子数）；$L = |\sum m_l|$。
  \item 不足半满时 $J = |L - S|$，超过半满时 $L + S$；
        恰好半满时 $L = 0$ 且 $J = S$。
\end{enumerate}
\end{method}

\begin{example}[基态项]\label{ex:b3:many-electron-atoms:ground-terms}
% ledger: lev:Ti2.3P0, lev:Ni2.3P0, lev:Cr3.4P1_2
\ce{Ti^2+}，$d^2$：$m_l = 2, 1$ 两个轨道框已填，$S = 1$，$L = 3$，$J = 2$：
\termsym{3}{F}{2}。\ce{Cr^3+}，$d^3$：$S = \frac32$，$L = 2 + 1 + 0 = 3$，
\termsym{4}{F}{3/2}。\ce{Fe^2+}，$d^6$：$m_l = 2$ 处五个自旋向上、一个向下，
$S = 2$，$L = 2$，超过半满，\termsym{5}{D}{4}。\ce{Ni^2+}，$d^8$：
$S = 1$，$L = 3$，\termsym{3}{F}{4}。所得结果均与原子光谱
测得的基态能级一致。
\end{example}

\begin{definition}[自旋–轨道耦合，精细结构]\label{def:b3:many-electron-atoms:spin-orbit}
称\emph{自旋–轨道耦合}\index{自旋–轨道耦合}为
电子自旋磁矩与其轨道运动磁矩之间的相互作用，
对某一项记为 $A\,\hat{\mathbf L}\cdot\hat{\mathbf S}/\hbar^2$。它
将一个项按不同的 $J$ 分裂为各能级，该项的\emph{精细
结构}\index{精细结构}。
\end{definition}

\begin{proposition}[朗德间隔定则]\label{prop:b3:many-electron-atoms:lande}
在同一项内，$E(J) = E_0 + \frac{A}{2}[J(J+1) - L(L+1) - S(S+1)]$，于是
$E(J) - E(J - 1) = AJ$。$A > 0$（正常）对应不足半满的壳层，
$A < 0$（倒置）对应超过半满的壳层。
\end{proposition}

\begin{proof}
由 $\hat{\mathbf J} = \hat{\mathbf L} + \hat{\mathbf S}$ 得 $\hat{\mathbf J}^2 =
\hat{\mathbf L}^2 + \hat{\mathbf S}^2 + 2\hat{\mathbf L}\cdot\hat{\mathbf S}$，于是
$\hat{\mathbf L}\cdot\hat{\mathbf S} = \frac12(\hat{\mathbf J}^2 - \hat{\mathbf L}^2 -
\hat{\mathbf S}^2)$，其在给定 $J$、$L$、$S$ 的态中的本征值为
$\frac{\hbar^2}{2}[J(J+1) - L(L+1) - S(S+1)]$。$J$ 与
$J - 1$ 的能级差为 $\frac A2[J(J+1) - (J-1)J] = AJ$。$A$ 的符号在本阶段不加证明地承认：超过半满的
壳层行为相当于相应数目的正
空穴。
\end{proof}

\begin{example}[检验间隔定则]\label{ex:b3:many-electron-atoms:lande-test}
% ledger: lev:C.3P1, lev:C.3P2, lev:O.3P1, lev:O.3P0
对于 ${}^3$P，该定则预言间隔之比为 $J = 2 : 1$，即 2。
碳：$(43.4 - 16.4)/16.4 = 1.65$。氧（$2p^4$，倒置，\termsym{3}{P}{2}
最低，其次 $J = 1$ 位于 \qty{158.3}{cm^{-1}}，$J = 0$ 位于 \qty{227.0}{cm^{-1}}）：
$158.3/68.7 = 2.30$。该定则在 20\,\% 以内成立：罗素–桑德斯耦合
是对轻原子的良好描述，但并非精确描述。
\end{example}

\begin{omfigure}
\begin{tikzpicture}[font=\small]
  \draw[omstate] (0,0) -- (1.6,0) node[midway, below] {$3s\ {}^2$S$_{1/2}$};
  \draw[omstate] (0,3.0) -- (1.6,3.0);
  \node[anchor=east] at (-0.1,3.0) {$3p$};
  \draw[omstate] (2.8,2.75) -- (4.4,2.75) node[right] {${}^2$P$_{1/2}$};
  \draw[omstate] (2.8,3.25) -- (4.4,3.25) node[right] {${}^2$P$_{3/2}$};
  \draw[densely dotted, gray] (1.6,3.0) -- (2.8,2.75) (1.6,3.0) -- (2.8,3.25);
  \draw[omfluo] (3.3,2.73) -- (3.3,0.02) node[pos=0.55, left] {D$_1$};
  \draw[omfluo] (3.9,3.23) -- (3.9,0.02) node[pos=0.55, right] {D$_2$};
  \draw[omstate] (2.8,0) -- (4.4,0);
  \draw[densely dotted, gray] (1.6,0) -- (2.8,0);
  \draw[decorate, decoration={brace, amplitude=3pt}] (5.5,3.27) -- (5.5,2.73)
     node[midway, right=3pt, align=left] {\footnotesize \qty{17.2}{cm^{-1}}\\[-1pt]\footnotesize （未按比例）};
\end{tikzpicture}

{\small 钠的 D 线。$3p$ 能级因自旋–轨道耦合
而分裂为 ${}^2$P$_{1/2}$ 与 ${}^2$P$_{3/2}$；向基态 $3s$ 能级发射
两条黄色谱线：D$_1$ 位于 \qty{589.76}{nm}，D$_2$ 位于 \qty{589.16}{nm}（真空
波长）。}
\end{omfigure}

\begin{example}[钠双线]\label{ex:b3:many-electron-atoms:sodium}
% ledger: lev:Na.3p1_2, lev:Na.3p3_2
钠的 $3p$ 能级位于 16956.17 与 \qty{16973.37}{cm^{-1}}：
D 线的真空波长为 $10^7/16956.17 = \qty{589.76}{nm}$ 与 \qty{589.16}{nm}，
二者相隔 \qty{17.20}{cm^{-1}}。对于 ${}^2$P，$E(\frac32) - E(\frac12) =
\frac32A$，故 $A = \qty{11.47}{cm^{-1}}$。
\end{example}

\begin{proposition}[原子的选择定则]\label{prop:b3:many-electron-atoms:selection-atoms}
在罗素–桑德斯耦合中，电偶极跃迁要求
$\Delta S = 0$，$\Delta L = 0, \pm1$，$\Delta J = 0, \pm1$（但不包括 $J = 0 \to J
= 0$），且宇称改变：对于单电子跃迁，$\Delta l = \pm1$。
\end{proposition}
\begin{proof}[说明]
$\Delta S = 0$ 已在\cref{ch:b3:electronic-spectroscopy}中证明：偶极
算符不作用于自旋。其余结论由偶极算符的矢量性质
得出，在本阶段不加证明地承认。
\end{proof}

\begin{omfigure}
\begin{tikzpicture}
\begin{axis}[width=0.62\linewidth, height=6.2cm, xmin=-0.5, xmax=3.5,
  ymin=158, ymax=173, axis y line=left, axis x line=none,
  ylabel={能量 / $10^3$ cm$^{-1}$}, unbounded coords=jump, clip=false,
  tick label style={font=\small}, label style={font=\small}]
\addplot[omstate] table[x=x, y=E] {figdata/out/bachelor-3-atomic-levels-helium.dat};
\node[font=\small, anchor=west] at (axis cs:1.05,166.277) {$1s2s\ {}^1$S};
\node[font=\small, anchor=west] at (axis cs:1.05,171.135) {$1s2p\ {}^1$P};
\node[font=\small, anchor=west] at (axis cs:3.05,159.856) {$1s2s\ {}^3$S};
\node[font=\small, anchor=west] at (axis cs:3.05,169.087) {$1s2p\ {}^3$P};
\node[font=\small, anchor=south] at (axis cs:0.5,172.3) {单重态};
\node[font=\small, anchor=south] at (axis cs:2.5,172.3) {三重态};
\draw[omfluo] (axis cs:0.5,171.0) -- (axis cs:0.5,166.42) node[pos=0.5, left, font=\scriptsize] {\qty{2059}{nm}};
\draw[omfluo] (axis cs:2.5,168.95) -- (axis cs:2.5,160.0) node[pos=0.5, left, font=\scriptsize] {\qty{1083}{nm}};
\draw[densely dotted, gray] (axis cs:-0.3,159.856) -- (axis cs:2,159.856);
\draw[<->, omThm] (axis cs:-0.2,159.856) -- (axis cs:-0.2,166.277) node[pos=0.5, right, font=\scriptsize] {$2K$};
\end{axis}
\end{tikzpicture}

{\small 两种“氦”。单重态与三重态能级，分属组态
$1s2s$ 与 $1s2p$ 的能级，取其在基态以上的实测能量（
${}^3$P 项取其各能级的平均值）。连线仅连接同一
族的态；每个三重态均比其单重态低 $2K$。}
\end{omfigure}

\begin{history}[两种“氦”与不相容原理]
氦于 1868 年因太阳而得名，1895 年在地球上由
拉姆齐发现。其两个谱线系令光谱学家困惑了三十
年。1925 年，泡利提出：任意两个电子的量子数不能全部相同，
同年，乌伦贝克与
古德斯密特提出了电子自旋；1926 年，海森堡证明
仅凭波函数的对称性即可将氦分裂为单重态与
三重态两族。
\end{history}

\begin{inthelab}[分辨钠双线]
将钠灯置于光栅光谱仪的狭缝前。若用
每毫米 600 条刻线的光栅，则两条谱线的一级衍射角
相差约 \num{0.02}\textdegree：优质光谱仪能将二者分开，学生用的
棱镜则不能。灯工作时温度很高，须待其
冷却后方可操作。
\end{inthelab}

\section{练习}

\begin{exercise}[$\star$]\label{exo:b3:many-electron-atoms:1}
写出锂的基态 $1s^22s$ 的斯莱特行列式，其中
$2s$ 电子的自旋为 $\alpha$。为何
$1s^3$ 没有对应的行列式？
\end{exercise}

\begin{exercise}[$\star$]\label{exo:b3:many-electron-atoms:2}
计算组态 $p^2$、$p^3$ 与 $d^2$ 的行列式个数。
$p^3$ 的谱项为 ${}^4$S、${}^2$D、${}^2$P，$d^2$ 的谱项为 ${}^3$F、${}^3$P、${}^1$G、
${}^1$D、${}^1$S：请核对两者计数。
\end{exercise}

\begin{exercise}[$\star$]\label{exo:b3:many-electron-atoms:3}
给出 N、O、F、\ce{Ti^2+} 与 \ce{Fe^3+}（$d^5$）的基态谱项与能级。
\end{exercise}

\begin{exercise}[$\star$]\label{exo:b3:many-electron-atoms:4}
列出谱项 ${}^2$D 与 ${}^3$P 的各能级及其简并度 $2J+1$，
并验证各简并度之和为 $(2L+1)(2S+1)$。
\end{exercise}

\begin{exercise}[$\star\star$]\label{exo:b3:many-electron-atoms:5}
碳的 ${}^3$P 能级位于 0、16.4 与 \qty{43.4}{cm^{-1}}。计算
各间隔之比，并由每个间隔求自旋–轨道耦合常数 $A$。
两个 $A$ 值的差异说明了什么？
\end{exercise}

\begin{exercise}[$\star\star$]\label{exo:b3:many-electron-atoms:6}
钠的 $3p$ 能级位于 16956.17 与 \qty{16973.37}{cm^{-1}}，计算
D 线的真空波长、它们的间隔（以 \unit{cm^{-1}} 与
meV 表示），以及 $A$ 值（$3p$ 谱项）。
\end{exercise}

\begin{exercise}[$\star\star$]\label{exo:b3:many-electron-atoms:7}
钠的下列跃迁中，哪些在发射中是允许的：$3p \to 3s$，
$3d \to 3s$、$3d \to 3p$、$4s \to 3s$、$4s \to 3p$？逐一说明理由。
\end{exercise}

\begin{exercise}[$\star\star$]\label{exo:b3:many-electron-atoms:8}
% ledger: lev:He.2p3P2, lev:He.2p3P1, lev:He.2p3P0, lev:He.2p1P1
氦的 $1s2p$ 能级为 \termsym{3}{P}{2}、\termsym{3}{P}{1}、
\termsym{3}{P}{0} 分别位于 169086.76、169086.84、\qty{169087.83}{cm^{-1}}，
\termsym{1}{P}{1} 位于 \qty{171134.89}{cm^{-1}}。
计算 ${}^3$P 谱项的平均能量，再求以 eV 计的 $K(1s,2p)$。
与 $K(1s,2s) = \qty{0.398}{eV}$ 相比，并解释其差异。
\end{exercise}

\begin{exercise}[$\star\star$]\label{exo:b3:many-electron-atoms:9}
用本章方法求 $d^2$ 的谱项，从
最大的 $M_L$ 出发（可以不写出完整表格，但须说明每个谱项的理由）。
\end{exercise}

\begin{exercise}[$\star\star\star$]\label{exo:b3:many-electron-atoms:10}
证明三重态的空间函数 $\frac{1}{\sqrt2}[1s(1)2s(2) -
2s(1)1s(2)]$ 在 $\mathbf r_1 = \mathbf r_2$ 时为零，单重态
函数则不为零。并将此与 $E_+ - E_-$ 的符号联系起来。
\end{exercise}

\begin{exercise}[$\star\star\star$]\label{exo:b3:many-electron-atoms:11}
% ledger: lev:Ni2.3F3, lev:Ni2.3F2
对 \ce{Ni^2+}（$d^8$），${}^3$F 能级位于 0（$J = 4$）、1360.7（$J = 3$）
及 \qty{2269.6}{cm^{-1}}（$J = 2$）。解释次序，计算
间隔之比，并与朗德的预言比较。由较大的
间隔求 $A$。
\end{exercise}

\begin{exercise}[$\star\star\star$]\label{exo:b3:many-electron-atoms:12}
证明 $\hat{\mathbf L}\cdot\hat{\mathbf S}$ 的期望值在
一个谱项的全部 $(2L+1)(2S+1)$ 个态上求和为零，从而自旋–轨道
耦合不改变能级的加权平均。在
碳的 ${}^3$P 上验证这一点。
\end{exercise}

\section{问题：两种氦}

\begin{problem}[{周末问题——氦的单重态与三重态两族：
行列式、过渡金属离子的谱项、区分两族的选择定则，
以及交换积分的测量}]
\label{pb:b3:many-electron-atoms:1}
% ledger: lev:He.2s3S1, lev:He.2s1S0, lev:He.2p3P2, lev:He.2p3P1, lev:He.2p3P0, lev:He.2p1P1, lev:Ti2.3F3, lev:Ti2.3F4, lev:Ti2.3P0, lev:Ti2.3P1, lev:Ti2.3P2
数据（NIST，以 \unit{cm^{-1}} 计，相对于各物种的基态）。氦：
$1s2s$ \termsym{3}{S}{1} 159855.97，\termsym{1}{S}{0} 166277.44；$1s2p$
\termsym{3}{P}{2,1,0} 169086.76、169086.84、169087.83，\termsym{1}{P}{1}
171134.89。\ce{Ti^2+}：\termsym{3}{F}{2,3,4} 为 0、184.9、420.4；\termsym{3}{P}{0,1,2}
10538.4、10603.6、10721.2。$\qty{1}{eV} = \qty{8065.54}{cm^{-1}}$。

\medskip\noindent\textbf{第一部分——自旋轨道与行列式。}
\begin{enumerate}
  \item 列出组态 $1s2s$ 可用的四个自旋轨道。
  \item 写出一个电子在 $1s$、另一个电子在
        $2s$ 中的四个行列式。
  \item 证明 $|1s\alpha\,2s\alpha|$ 是反对称空间函数与自旋函数
        $\alpha(1)\alpha(2)$ 的乘积。
  \item 组合 $|1s\alpha\,2s\beta|$ 与 $|1s\beta\,2s\alpha|$，
        得到同一空间函数的 $M_S = 0$ 分量及第四个态。
  \item 哪种空间函数（对称或反对称）对应
        $S = 0$？对应 $S = 1$？
  \item 为什么称它们为单重态与三重态？
\end{enumerate}

\medskip\noindent\textbf{第二部分——\ce{Ti^2+} 的谱项。}
\begin{enumerate}[resume]
  \item 给出 \ce{Ti^2+} 的组态及其
        行列式的个数。
  \item 用洪特定则求其基态谱项与能级，对照
        给出的数据进行核对。
  \item 其余谱项为 ${}^3$P、${}^1$G、${}^1$D 与 ${}^1$S，核对
        态的数目。
  \item ${}^3$F 的精细结构是正常的还是倒置的？检验朗德
        间隔定则。
  \item 计算 ${}^3$F 与 ${}^3$P 谱项的加权平均能量。
  \item 二者的间隔为 $15B$，其中 $B$ 为该离子的拉卡参数
        （见\cref{ch:b3:complex-spectra-magnetism}）。计算 $B$。
\end{enumerate}

\medskip\noindent\textbf{第三部分——为何分为两族。}
\begin{enumerate}[resume]
  \item 哪条选择定则禁止单重态与
        三重态之间的跃迁？
  \item 最低的三重态 $1s2s$ \termsym{3}{S}{1} 不能向基态
        $1s^2$ \termsym{1}{S}{0} 发射。给出两条理由，
        并说明这类态称为什么。
  \item 计算谱线 $1s2p\ {}^3$P $\to 1s2s\ {}^3$S 的波长（取
        ${}^3$P 能级的平均值）。
  \item 计算 $1s2p\ {}^1$P $\to 1s2s\ {}^1$S 的波长。
  \item 计算 $1s2p\ {}^1$P $\to 1s^2\ {}^1$S 的波长。它位于
        光谱的哪个区域？
  \item 解释十九世纪的光谱学家为何在看到这些谱线时，
        相信存在两种不同的气体。
\end{enumerate}

\medskip\noindent\textbf{第四部分——交换积分的测量。}
\begin{enumerate}[resume]
  \item 写出 $1s2s$ 单重态与三重态的一级能量，
        用 $E_{1s}$、$E_{2s}$、$J$ 与 $K$ 表示。
  \item 哪一个更低？从电子的位置来看，原因是什么？
  \item 计算 $1s2s$ 单重态与三重态之间的能隙，以 \unit{cm^{-1}} 与 eV 计。
  \item 以 eV 计算 $K(1s,2p)$，所用的是 $1s2p$ 各能级。
  \item 为什么 $K(1s,2s)$ 大于 $K(1s,2p)$？
  \item 氦的激发态是否符合洪特第一定则？
  \item 给出结果：氦的交换积分 $K(1s,2s)$，
        以 eV 计。
\end{enumerate}
\end{problem}
